\section*{Capítulo \ref{ch:b3:electrons-in-solids} --- Electrones en los sólidos}

\begin{solution}{exo:b3:electrons-in-solids:1}
(a) $\tau = m_{\text{e}}/ne^2\rho = \qty{2.5e-14}{s}$. (b) $\ell =
v_{\text{F}}\tau \approx \qty{39}{nm}$. (c) Unas 110 \omterm{def:b3:crystalline-solids:lattice}{constantes
de red}. (d) Un electrón que navega ante un centenar de iones sin
dispersarse no puede estar rebotando en los propios iones --- la
\omterm{def:b3:crystalline-solids:lattice}{red} ha de ser transparente para él, exactamente lo que el \omterm{thm:b3:electrons-in-solids:bloch}{teorema de Bloch}
(\cref{thm:b3:electrons-in-solids:bloch}) demuestra después.
\end{solution}

\begin{solution}{exo:b3:electrons-in-solids:2}
(a) $v_{\text{d}} = I/neA = \qty{7.3e-4}{m/s}$ --- menos de un
milímetro por segundo. (b) Unos 23 minutos. (c) El \emph{campo}
se establece a lo largo del hilo casi a la velocidad de la luz y pone
a la deriva de golpe a todo el mar de electrones: los que marchan son lentos,
la orden de marchar no. (d) $v_{\text{d}}/v_{\text{F}} \sim
10^{-9}$: la conducción es un sesgo imperceptible sobre un violento
movimiento cuántico, no un flujo clásico apacible.
\end{solution}

\begin{solution}{exo:b3:electrons-in-solids:3}
(a) Banda semillena: \omterm{def:b3:electrons-in-solids:classes}{metal}, $\dd\rho/\dd T > 0$. (b) Dos
electrones llenarían su banda exactamente --- el magnesio conduce
porque su banda llena \emph{se solapa} con la siguiente, vacía: \omterm{def:b3:electrons-in-solids:classes}{metal},
$\dd\rho/\dd T > 0$. (c) Banda llena y \omterm{thm:b3:electrons-in-solids:bloch}{banda prohibida} de \qty{5.5}{eV}: \omterm{def:b3:electrons-in-solids:classes}{aislante}
(formalmente $\dd\rho/\dd T < 0$, pero sin apenas portadores
que contar). (d) \omterm{def:b3:electrons-in-solids:classes}{Semiconductor}: $\dd\rho/\dd T < 0$, el
cambio exponencial de signo que delata la clase.
\end{solution}

\begin{solution}{exo:b3:electrons-in-solids:4}
(a) $\lambda = hc/E_{\text{g}} = \qty{1.1}{\micro m}$: el silicio es
transparente en el infrarrojo más allá de ahí --- y por eso las cámaras de
infrarrojos pueden llevar lentes de silicio. (b) \qty{225}{nm}, en el
ultravioleta: ningún fotón visible puede absorberse, así que el diamante es
transparente como el agua. (c) \qty{365}{nm}: una \omterm{thm:b3:electrons-in-solids:bloch}{banda prohibida} lo bastante ancha para el azul
(\qty{2.8}{eV}) existe en los nitruros y casi en ningún otro sitio
práctico --- el LED azul esperó décadas al crecimiento de materiales.
(d) Su \omterm{thm:b3:electrons-in-solids:bloch}{banda prohibida} fija su fotón: sobreexcitarlo añade calor, no anchura
de banda, y el diodo se cuece a oscuras --- el color de un LED es una
constante del material, no un ajuste de brillo.
\end{solution}

\begin{solution}{exo:b3:electrons-in-solids:5}
(a) $\kappa/\sigma T = 400/(5.9\times10^7\times300) =
\qty{2.3e-8}{W.\ohm/K^2}$. (b) A menos de un diez por ciento de $L$. (c)
Los mismos electrones de la superficie de Fermi llevan tanto la carga como el
calor, así que su tiempo de dispersión se cancela en el cociente. (d) La
sartén de acero conduce el calor a la comida porque sus electrones se mueven;
el mango de madera, \omterm{def:b3:electrons-in-solids:classes}{aislante} en ambos sentidos, los mantiene ---
y al calor con ellos --- fuera de su mano.
\end{solution}

\begin{solution}{exo:b3:electrons-in-solids:6}
(a) Un cincuenta por ciento de más --- flagrantemente ausente del experimento.
(b) $T/T_{\text{F}} \approx 0.004$: el término electrónico queda
estrangulado por debajo del uno por ciento. (c) El $T^3$ de la \omterm{def:b3:crystalline-solids:lattice}{red} se desploma
más deprisa que la $T$ de los electrones: por debajo de unos pocos kelvin, los
electrones ``despreciables'' son lo único que queda. (d) $C/T =
\gamma + \beta T^2$: la ordenada en el origen $\gamma$ pesa los electrones y
la pendiente $\beta$ la \omterm{def:b3:crystalline-solids:lattice}{red} --- una gráfica, los dos inquilinos.
\end{solution}

\begin{solution}{exo:b3:electrons-in-solids:7}
(a) $k_{\text{F}} = (3\pi^2n)^{1/3} = \qty{1.36e10}{m^{-1}}$:
$\lambda_{\text{F}} = \qty{4.6}{\angstrom}$. (b) $2a =
\qty{7.2}{\angstrom}$: el mismo orden --- la superficie de Fermi del cobre
llega hasta cerca del borde de zona, y la física que abre la \omterm{thm:b3:electrons-in-solids:bloch}{banda prohibida}
ocurre \emph{en} las energías que importan. (c) Una onda estacionaria
amontona su carga sobre los iones positivos (menor energía
electrostática) y la otra entre ellos (mayor): la misma longitud de onda, dos
energías --- la \omterm{thm:b3:electrons-in-solids:bloch}{banda prohibida}. (d) El recorrido libre de 110 \omterm{def:b3:crystalline-solids:lattice}{constantes de red}:
las ondas de Bloch no se dispersan en una \omterm{def:b3:crystalline-solids:lattice}{red} perfecta, de modo que solo quedan
las vibraciones y las impurezas para ofrecer resistencia.
\end{solution}

\begin{solution}{exo:b3:electrons-in-solids:8}
(a) $E_{\text{g}}/2k_{\text{B}} = \qty{6500}{K}$:
$\exp[6500(1/300 - 1/400)] \approx 230$. (b) La resistencia del
termistor se desploma exponencialmente --- una curva abrupta y sensible;
la del platino sube con suavidad y linealmente (más fonones). Uno es un
termómetro por recuento de portadores y el otro por dispersión. (c)
Los portadores nacen en \emph{parejas} de electrón y \omterm{def:b3:electrons-in-solids:doping}{hueco}, y el
equilibrio $np = n_{\text{i}}^2$ reparte el coste de la \omterm{thm:b3:electrons-in-solids:bloch}{banda prohibida} entre
los dos --- de ahí el $E_{\text{g}}/2$. (d)
$n_{\text{i}}$ alcanza $\qty{5e21}{m^{-3}}$ cerca de \qty{760}{K}:
el \omterm{def:b3:electrons-in-solids:doping}{dopado} se ahoga y el circuito olvida su diseño. Los
dispositivos reales se rinden antes --- sus fugas inversas, que se duplican cada
diez kelvin, se portan mal primero.
\end{solution}

\begin{solution}{exo:b3:electrons-in-solids:9}
(a) $\qty{5e22}{m^{-3}}$ --- cinco millones de veces $n_{\text{i}}$.
(b) El mismo factor $\sim5\times10^6$: una mota por cada millón de
átomos se adueña por completo de la conductividad. (c) Porque unas ppm accidentales
harían lo mismo sin que nadie las invite: solo un \omterm{def:b3:crystalline-solids:lattice}{cristal} puro hasta ppb tiene una
conductividad que pertenece al diseñador. (d) $d \sim n^{-1/3}
\approx \qty{27}{nm}$, diez veces la órbita de \qty{2.4}{nm}: los \omterm{def:b3:electrons-in-solids:doping}{donantes}
son hidrógenos aislados; súbase el \omterm{def:b3:electrons-in-solids:doping}{dopado} cien veces y las
órbitas se tocan --- una banda de impurezas y, al final, un \omterm{def:b3:electrons-in-solids:classes}{metal}.
\end{solution}

\begin{solution}{exo:b3:electrons-in-solids:10}
(a) $E = 13.6\times0.26/11.7^2 \approx \qty{26}{meV}$ (el
valor medido del fósforo, \qty{45}{meV}, mantiene honrada la
escala). (b) Comparable a $k_{\text{B}}T = \qty{26}{meV}$:
esencialmente todos los \omterm{def:b3:electrons-in-solids:doping}{donantes} están ionizados a temperatura ambiente --- la
premisa de la \cref{def:b3:electrons-in-solids:doping}. (c) $a =
\qty{0.53}{\angstrom}\times\varepsilon_{\text{r}}/(m^*/
m_{\text{e}}) \approx \qty{2.4}{nm}$. (d) Una órbita que abarca
decenas de celdas ve el \omterm{def:b3:crystalline-solids:lattice}{cristal} como un medio liso --- que es
precisamente la condición bajo la cual una $\varepsilon_{\text{r}}$
de volumen y una $m^*$ promediada sobre la banda son legítimas:
la aproximación se certifica a sí misma.
\end{solution}

\begin{solution}{exo:b3:electrons-in-solids:11}
(a) $V_{\text{bi}} = 0.0259\ln(10^{12}) \approx \qty{0.71}{V}$.
(b) $eV/k_{\text{B}}T = 19.3$: razón $\eu^{19.3} \approx
2\times10^{8}$ --- el diodo es una calle de sentido único con ocho
cifras. (c) $I_0 \propto n_{\text{i}}^2 \propto
\eu^{-E_{\text{g}}/k_{\text{B}}T}$: la misma exponencial que en el
\cref{exo:b3:electrons-in-solids:8}, de ahí la
duplicación de la regla del pulgar. (d) Los cuatro diodos forman un rombo: en cada
semiciclo, la pareja que quede en directa dirige la corriente
por la carga en el \emph{mismo} sentido --- entra alterna y sale continua
con baches.
\end{solution}

\begin{solution}{exo:b3:electrons-in-solids:12}
(a) Estréchese la \omterm{thm:b3:electrons-in-solids:bloch}{banda prohibida} y más fotones la superarán, pero cada uno entrega
solo la pequeña energía de la banda; ensánchese y cada fotón paga más
pero menos cumplen el requisito: la cosecha $\times$ tensión tiene su máximo en medio.
(b) Justo por debajo del óptimo de \qty{1.3}{eV}: el techo ideal del
silicio es del $\sim30\,\%$ --- lo bastante cerca para que gane su baratura.
(c) La célula superior de banda ancha se queda el azul a tensión alta y pasa
el rojo a la célula de banda estrecha de debajo: cada fotón se cosecha
cerca de su propia energía, la termalización se encoge y el techo de la
pila trepa hacia los cuarenta. (d) El calor sube $I_0$
exponencialmente y arrastra a la baja la tensión de circuito abierto
$\sim(k_{\text{B}}T/e)\ln(I_{\text{L}}/I_0)$; y la propia \omterm{thm:b3:electrons-in-solids:bloch}{banda prohibida}
se estrecha un poco, y cambia tensión por una corriente que no la
recupera del todo.
\end{solution}

\begin{solution}{pb:b3:electrons-in-solids:1}
\textbf{1.} $\lambda_{\text{max}} \approx \qty{500}{nm}$:
alrededor de \qty{2.5}{eV}.
\textbf{2.} $hc/E_{\text{g}} \approx \qty{1.1}{\micro m}$.
\textbf{3.} Atraviesa las células (no hay estado final que la
absorba) y acaba como calor en el respaldo --- calentando el tejado, no
los cables.
\textbf{4.} $\qty{1.12}{eV}$ sobreviven como el par; el exceso de
\qty{1.4}{eV} se escurre hacia los fonones en picosegundos --- más deprisa
de lo que ningún circuito podría interceptar.
\textbf{5.} La transparencia por debajo de la banda ($\sim20\,\%$ de la potencia) y
la termalización por encima ($\sim30\,\%$): al espectro le gravan
la mitad antes de que la unión vea un solo electrón.
\textbf{6.} La absorción es un salto cuántico de un solo fotón que necesita un
estado final real; a media \omterm{thm:b3:electrons-in-solids:bloch}{banda prohibida} no hay ninguno donde hacer escala, y
los sucesos de dos fotones son despreciables a las densidades de fotones de la luz solar.
\textbf{7.} Los pares creados en la \omterm{prop:b3:electrons-in-solids:junction}{zona de deplexión} o cerca de ella sienten el
campo de contacto: los electrones son barridos al lado n y los \omterm{def:b3:electrons-in-solids:doping}{huecos} al
lado p, y el campo los separa más deprisa de lo que pueden encontrarse
para recombinarse --- la carga se acumula hasta que un circuito
externo la alivia como corriente.
\textbf{8.} $V_{\text{bi}} = 0.0259\ln(10^{12}) \approx
\qty{0.71}{V}$.
\textbf{9.} $60\times\qty{0.6}{V} = \qty{36}{V}$.
\textbf{10.} $I = 400/36 \approx \qty{11}{A}$.
\textbf{11.} En todo el panel, la luz por encima de la \omterm{thm:b3:electrons-in-solids:bloch}{banda prohibida} son
$\sim\qty{1400}{W}$ a $\sim\qty{2}{eV}$ cada fotón: $4\times10^{21}$
fotones por segundo --- pero las sesenta células están en \emph{serie},
así que los \qty{11}{A} (un flujo de electrones $I/e \approx
7\times10^{19}\,\unit{s^{-1}}$) pasan por cada célula, cuya
ración propia de fotones es $4\times10^{21}/60 \approx
7\times10^{19}$ por segundo: casi todo fotón absorbido rinde
un electrón recogido. La eficiencia cuántica es excelente; las
pérdidas son energéticas, no numéricas.
\textbf{12.} Tras el 50\,\% del espectro, la célula paga el
déficit de tensión ($0.6/1.12 \approx 54\,\%$) y unos pocos puntos de
reflexión y recombinación: $0.5\times0.54 \approx 27\,\%$,
recortado hasta el 20\,\% nominal.
\textbf{13.} La fuga $I_0$: la tensión de circuito abierto
$(k_{\text{B}}T/e)\ln(I_{\text{L}}/I_0)$ se detiene donde
se equilibran la fotocorriente y la fuga del diodo --- alrededor de \qty{0.6}{V}, muy
lejos de la \omterm{thm:b3:electrons-in-solids:bloch}{banda prohibida}.
\textbf{14.} $\sin30^\circ = 0.5$: la mitad de lo nominal de mediodía solo por
geometría --- antes de las nubes.
\textbf{15.} $\Delta T = \qty{40}{K}$: un $-16\,\%$, que deja
$\approx\qty{335}{W}$ --- los días más soleados no son los mejores
por vatio.
\textbf{16.} El calor infla $I_0$ exponencialmente, y la
tensión de circuito abierto --- reprimenda de un logaritmo --- cae una
fracción de punto porcentual por kelvin; la banda ligeramente encogida
remata la faena.
\textbf{17.} El cableado en serie impone una corriente común: la célula
sombreada, que no genera ninguna, la meten en polarización inversa sus
cincuenta y nueve compañeras y disipa la potencia de estas como un punto caliente
--- la cadena se estrangula hasta la célula más débil.
\textbf{18.} El diodo de paso da a la corriente un camino en directa
alrededor de la subcadena sombreada, y sacrifica su tensión en vez de
la salida del panel entero.
\textbf{19.} $\qty{400}{W}\times0.15\times\qty{8760}{h} \approx
\qty{530}{kWh}$ al año.
\textbf{20.} $\approx\qty{130}{euros}$ al año: amortización en
unos $300/130 \approx 2.5$ años, y luego dos décadas de beneficio.
\textbf{21.} $500/530 \approx$ un año: el panel devuelve su
energía de fabricación casi tan deprisa como su precio.
\textbf{22.} La negrura no es una elección, sino una \omterm{thm:b3:electrons-in-solids:bloch}{estructura de bandas}:
la absorción necesita un estado vacío una energía de fotón por encima de uno lleno,
y por debajo de la \omterm{thm:b3:electrons-in-solids:bloch}{banda prohibida} el silicio sencillamente no tiene ninguno que ofrecer ---
la pintura no puede añadir estados.
\textbf{23.} Apílese encima una célula de banda ancha que se quede el azul a
tensión alta y deje pasar el rojo hacia abajo: el tándem bate el
techo de una sola unión. La pega: la pila en serie debe casar
las corrientes (y los precios) entre capas.
\textbf{24.} Todo lo que acorte $\tau$ y envenene las uniones:
los defectos creados por el ultravioleta y las impurezas que se difunden dentro dispersan y atrapan
portadores, la humedad corroe los contactos y los puntos calientes envejecen los
diodos --- la lenta entropía de un \omterm{def:b3:crystalline-solids:lattice}{cristal} al que se pide sentarse al sol
treinta años.
\textbf{25.} La \omterm{thm:b3:electrons-in-solids:bloch}{banda prohibida} cosecha la mitad del sol; la unión convierte
los pares en \qty{36}{V} de corriente ordenada; las cadenas en serie,
la sombra y el calor del verano se llevan su parte; y con
$\sim\qty{530}{kWh}$ al año contra una amortización energética de un año y
otra monetaria de tres, el árbitro dictamina: póngase en el tejado.
\end{solution}
